\section{What the edits explain}
\label{sec:explain}

An edit can improve answers without explaining why the original model failed. We therefore distinguish
\emph{repair}, a beneficial change in behaviour; \emph{pathway evidence}, an intervention linking the edit's
effect to a named component; and a \emph{mechanistic account}, a description that also predicts the unedited
model's behaviour. The repair result above concerns the layer-35 edit. The pathway evaluation selected a
different candidate: the two-head edit, at dose 1.5, with the layer-33 head fixed in advance as the primary
component.

\paragraph{A pathway through layer 33.}
Both pathway tests passed after Holm correction across the three branch-M tests ($p < 0.001$).
Reverting the layer-33 component in the edited model reduced recovery on false abstentions by 0.293
(95\% CI 0.225 to 0.368). The edit also recovered 0.717 more than the mean of five size-matched random edits,
B6. These tests address the effect of changing that component and the comparison with random perturbations.

Two activation interventions asked whether the effect could be transferred or blocked. Adding the edit's
change to the layer-33 output at the answer position in the unedited model recovered 0.808 of false
abstentions. Clamping that output back to its original value in the edited model removed only 0.193.
The layer-33 intervention was therefore sufficient, but not necessary, to release answers under these
conditions. The remaining recovery is consistent with another route, possibly through the layer-32 head;
that explanation was not directly tested. Without a random-head comparison, these results do not establish
that the two selected heads have a unique role.

\paragraph{The explanation was not established.}
The agent's account of what the component computes was not supported by the frozen tests. Its proposed
variable added no predictive value for false abstention, although the baseline predictor already had an
AUROC of 0.998 and left little room for improvement. Intervening on the variable also failed to establish a
different effect on answerable versus unanswerable questions relative to a matched logit bias: the
interaction was 0.0005 (95\% CI $-0.025$ to 0.027). Two activity predictions held, but the combined account
test did not pass.

This negative has an important qualification. The agent described the heads as an evidence-independent
``unknown'' prior. A null interaction is close to what that description would predict, so the frozen predictor
and interaction tests fit this account poorly. The result means that the account was not shown by the
pre-registered tests; it is not a refutation of that account.

At its pathway-test dose, the two-head edit was also not selective: on S it produced unsupported answers on
96\% of unanswerable rows. The pathway result identifies where an answer-release effect can be produced.
It does not explain the selected layer-35 repair.

\paragraph{Exploratory observations.}
On E, doubling the two heads changed 180 of 572 correct answers (31\%) into ``unknown'', within the agent's
prediction of 30--40\%. Removing the heads recovered most false abstentions on that split. These observations
motivated the bidirectional interpretation, but E was the split the agent studied and cannot independently
confirm it. The exploratory results are collected in \cref{app:e}.
